\chapter{纠缠、混合态与Bell实验}
\label{ch:entanglement}

\section{密度算符与混合态}

上一章的双光子实验统计的是两个输出端各收到一个光子的次数。若只测量其中一路，怎样描述这一路的量子态？
先看一个较简单的问题。设自旋制备装置每次以相等概率输出 $\ket{+z}$ 或 $\ket{-z}$，实验者没有记录每次选了哪一种。
沿 $z$ 轴测量时，正负读数各占一半，与每次都制备 $\ket{+x}$ 的情形相同。改沿 $x$ 轴测量，随机制备的束流仍是正负各半，$\ket{+x}$ 束却全给正读数。

一般地，设装置以概率 $p_r$ 制备 $\ket{\psi_r}$。对这一制备，粒子从分析器的 $\ket a$ 出口通过的概率为 $\abs{\braket{a}{\psi_r}}^2$。对各种制备取加权平均，便得到总通过率，
\begin{equation}
    \P(A = a) = \sum_r p_r\abs{\braket{a}{\psi_r}}^2
        = \mel{a}{\sum_rp_r\ketbra{\psi_r}{\psi_r}}{a}.
\end{equation}
右端夹在 $\bra a$ 和 $\ket a$ 之间的算符只取决于制备，称为密度算符 $\op\rho$。改变分析方向，只需换入相应的 $\ket a$。

把这两种制备写在同一组基$(\ket{+z},\ket{-z})$中，计数的差别就体现在矩阵元上。
\begin{equation}
    \op{\rho}_{\mixed} \longleftrightarrow \frac12\begin{pmatrix}1&0\\0&1\end{pmatrix},
    \qquad
    \op{\rho}_{+x} \longleftrightarrow \frac12\begin{pmatrix}1&1\\1&1\end{pmatrix}.
\end{equation}
这里的 $+$、$-$ 是 $z$ 基的行列标号。例如 $\rho_{+-} = \mel{+z}{\op{\rho}}{-z}$ 是一个复数。两种制备的对角元相同，沿 $z$ 轴测量都给出两端各半。对任意两态密度算符，
把$\bra{+x} = (\bra{+z} + \bra{-z})/\sqrt2$从左乘、$\ket{+x}$从右乘，
得到 $\P(+x) = (\rho_{++} + \rho_{--} + \rho_{+-} + \rho_{-+})/2$。
混合束的这个概率是$1/2$，纯$\ket{+x}$束为1。
用同一组基描述时，非对角元保留了在另一方向可以读出的相干信息。

\begin{definition}{密度算符与迹}{qm-density}
    密度算符(density operator) $\op{\rho}$是正的、迹为1的算符。
    在正交归一基 $\ket{a}$ 中，算符 $\op{A}$ 的迹是其矩阵对角元之和，
    \begin{equation}
        \tr \op{A} = \sum_a\mel{a}{\op{A}}{a}.
    \end{equation}
    迹直接以算符为输入，在有限维对所有算符有定义，在无限维可取迹类算符为定义域。迹与所选基底无关。对有限维算符，或一个迹类算符与一个有界算符，满足循环性质 $\tr(\op{A}\op{B}) = \tr(\op{B}\op{A})$。
    正性(positivity) $\op{\rho} \geq 0$指每个$\ket{u}$都满足$\mel{u}{\op{\rho}}{u} \geq 0$。
\end{definition}

若以经典概率$p_r$制备$\ket{\psi_r}$，忽略制备标签，定义密度算符
\begin{equation}
    \op{\rho} = \sum_rp_r\ketbra{\psi_r}{\psi_r},\qquad
    \op{\rho} \geq 0,\quad\tr\op{\rho} = 1.
\end{equation}
设 $\op\Pi_a$ 为可观测量 $\op A$ 对应于本征值 $a$ 的正交投影。先取有界的 $\op A$，平均读数和测得 $a$ 的概率分别为
\begin{subequations}
    \begin{align}
        \ave{\op{A}} &= \tr(\op{\rho}\op{A}),\\
        p_a &= \tr(\op{\rho}\op\Pi_a). \label{eq:density-statistics}
    \end{align}
\end{subequations}
投影 $\op\Pi_a$ 总是有界，所以概率公式也适用于无界可观测量。无界量的平均读数须按其谱概率分布求一阶矩（\ssecref{ss:qm-spectral-functions}）。
现在可以把\ssecref{ss:rabi-measurement}中由计数定义的布居写成密度矩阵元。在所选的正交归一基中，态 $\ket a$ 的布居为 $p_a = \mel{a}{\op\rho}{a}$。
测量若只分辨能量，同一简并能级的布居就是该能级内一组正交归一基矢的布居之和。上面两种制备在 $z$ 基中的布居同为 $(1/2,1/2)$，相干的差别在 $x$ 基中表现为不同的计数。
不同的制备系综(ensemble)可以给出同一个$\op{\rho}$，因而对系统的任何测量都有相同的结果分布。
纯态满足$\op{\rho}^2 = \op{\rho}$、$\tr\op{\rho}^2 = 1$，混合态(mixed state)有$\tr\op{\rho}^2 < 1$。

对随机制备的各个纯态分别求期望，再按制备概率加权，就可以验证密度算符的正性、迹归一以及平均读数公式，
\begin{subequations}
    \begin{alignat}{1}
        \mel{u}{\op{\rho}}{u} &= \sum_rp_r\abs{\braket{u}{\psi_r}}^2 \geq 0,\\
        \tr\op{\rho} &= \sum_rp_r\sum_a\abs{\braket{a}{\psi_r}}^2 = \sum_rp_r = 1,\\
        \tr(\op{\rho} \op{A}) &= \sum_rp_r\mel{\psi_r}{\op{A}}{\psi_r}.
    \end{alignat}
\end{subequations}
再将 $\op{\rho}$ 对角化，得到非负的本征值 $\lambda_j$。在本征基中有
\begin{subequations}\label{eq:qm-density-purity}
    \begin{gather}
        \op{\rho} = \sum_j\lambda_j\ketbra{j}{j},\qquad
        \sum_j\lambda_j = 1,\\
        \tr\op{\rho}^2 = \sum_j\lambda_j^2 \leq 1.
    \end{gather}
\end{subequations}
等号只有在一个$\lambda_j = 1$、其余为零时成立，
这时$\op{\rho}$就是一维投影。

两态系统的密度算符可以唯一写成
\begin{equation}
    \op{\rho} = \frac12(\one + \bm r \cdot \opbm{\sigma}),\qquad \abs{\bm r} \leq 1.
\end{equation}
本征值为$(1 \pm \abs{\bm r})/2$，故所有物理态都在 Bloch 球内或球面上，纯态恰在球面。

$\bm r = (r_x,r_y,r_z)$称为 Bloch 向量，其分量由自旋测量确定，
\begin{equation}
    r_i = \tr(\op{\rho}\op{\sigma}_i) = \ave{\op{\sigma}_i}.
\end{equation}
因为 $\one,\op{\sigma}_x,\op{\sigma}_y,\op{\sigma}_z$ 在实系数下张成两态系统的全部自伴算符，
迹为1固定了$\one$的系数为$1/2$。
其余三个系数由$\tr(\op{\sigma}_i\op{\sigma}_j) = 2\delta_{ij}$读出。
由于 $\bm r \cdot \opbm{\sigma}$ 的迹为零，平方又等于 $\abs{\bm r}^2\one$，它的本征值为 $\pm\abs{\bm r}$。加上单位算符再除以二，便得到 $\op\rho$ 的两个本征值。

对封闭系统，$\op{\rho}(t) = \op{U}(t)\op{\rho}(0)\op{U}\adjoint(t)$，求导便得von Neumann方程
\begin{equation}
    \i\hbar\dot{\op{\rho}} = [\op{H},\op{\rho}].
\end{equation}
幺正演化保持$\op{\rho}$的本征值，所以保持纯度(purity)和熵(entropy)。
当 Hamilton 算符不显含时间时，满足$[\op{H},\op{\rho}] = 0$的态是平稳态，可以是混合态；
能量基中的非对角元以相应能量差频率变化，简并态间的相干也可以保持不变。

\section{部分迹与约化密度算符}
上节的混合态来自舍去制备标签。设两台分析器分别测量联合系统的两端 $\subsysA,\subsysB$。
若只保留 $\subsysA$ 端的记录，将联合计数表中 $a$ 相同的各项相加，就得到边缘概率 $p_{\subsysA}(a)$。
要预言转动 $\subsysA$ 分析器后的计数，还需要找出对每一种测量方向都适用的局部密度算符。

在 $\subsysA,\subsysB$ 两端各选一组正交归一基，按先 $\subsysA$ 后 $\subsysB$ 的次序组成乘积基 $\ket{a,b} = \ket{a}_{\subsysA} \otimes \ket{b}_{\subsysB}$。任意联合纯态都可在这组基中展开，
\begin{equation}\label{eq:qm-bipartite-amplitudes}
    \ket{\Psi} = \sum_{a,b}C_{ab}\ket{a,b},\qquad C_{ab} = \braket{a,b}{\Psi}.
\end{equation}
把 $\subsysA$ 端的分析器转到 $\ket{u}$，仍忽略 $\subsysB$ 端的结果，则
\begin{equation}
    p_{\subsysA}(u) = \sum_b\Bigabs{\sum_a\braket{u}{a}C_{ab}}^2
    = \sum_{a,a'}\braket{u}{a}
        \Bigp{\sum_b C_{ab}C_{a'b}\conj}\braket{a'}{u}.
\end{equation}
固定$b$时，通向同一结果$(u,b)$的各项先相加振幅。
不同$b$代表互相正交的结果，忽略读数时再相加概率。括号内的系数是 $\subsysA$ 端约化密度算符的矩阵元，记这个算符为 $\op{\rho}_{\subsysA}$。
\begin{definition}[unbreakable]{部分迹与约化密度算符}{qm-partial-trace}
    对联合态 $\op{\rho}_{\subsysAB}$，取 $\subsysB$ 的任意一组正交归一基 $\ket{b}$，对 $\subsysB$ 作部分迹(partial trace)得到 $\subsysA$ 的约化密度算符
    \begin{equation}
        \op{\rho}_{\subsysA} = \tr_{\subsysB}\op{\rho}_{\subsysAB}
        = \sum_b(\one_{\subsysA} \otimes \bra{b})\op{\rho}_{\subsysAB}(\one_{\subsysA} \otimes \ket{b}).
    \end{equation}
    它对任意局部可观测量 $\op{O}_{\subsysA}$ 满足 $\tr_{\subsysA}(\op{\rho}_{\subsysA} \op{O}_{\subsysA}) = \ave{\op{O}_{\subsysA} \otimes \one_{\subsysB}}$。
\end{definition}

若
$\rho_{(a,b),(a',b')} = \mel{a,b}{\op{\rho}_{\subsysAB}}{a',b'}$，
部分迹的分量是
\begin{equation}
    (\op{\rho}_{\subsysA})_{aa'} = \sum_b\rho_{(a,b),(a',b)}.
\end{equation}
验证它给出全部局部测量只需展开迹：
\begin{equation}
    \tr_{\subsysAB}\bigb{\op{\rho}_{\subsysAB}(\op{O}_{\subsysA} \otimes \one_{\subsysB})}
    = \sum_{a,a',b}\rho_{(a,b),(a',b)}(\op{O}_{\subsysA})_{a'a}
    = \tr_{\subsysA}(\op{\rho}_{\subsysA}\op{O}_{\subsysA}).
\end{equation}
在$\subsysB$上更换正交基时，变换矩阵与其伴随相乘为单位矩阵，
故这个求和不变。

\begin{example}{双自旋态的单侧测量}{qm-marginal-counts}
    两端的自旋制备在 $\ket{\Psi} = (\ket{+z,+z} + \ket{-z,+x})/\sqrt2$，第一个标签对应 $\subsysA$ 端。两端都测 $z$ 时，按 $(+,+),(+,-),(-,+),(-,-)$ 排序的概率为 $1/2,0,1/4,1/4$。现在将 $\subsysA$ 端的分析器转向 $x$，求在该端测得正结果的概率。若只知道原来的四项计数，能否作出同一预言？
    \par\smallskip\textbf{解}\quad
    原先在 $\subsysA$ 端测得正、负结果的概率均为 $1/2$。
    用 $z$ 基展开 $\ket{+x}_{\subsysB}$。系数矩阵的行、列分别对应 $\subsysA,\subsysB$ 两端的基态，在这些基底中有
    \begin{subequations}
        \begin{align}
            C &= \begin{pmatrix}1/\sqrt2&0\\1/2&1/2\end{pmatrix},\\
            \op{\rho}_{\subsysA} &\longleftrightarrow CC\adjoint
            = \frac12\begin{pmatrix}1&1/\sqrt2\\1/\sqrt2&1\end{pmatrix}.
        \end{align}
    \end{subequations}
    因此 $p_{\subsysA}(+x) = (1 + 1/\sqrt2)/2$。
    \begin{notecaptionblock}
        \tikzfigure{reduced-spin-density-matrices}
        \captionof{figure}{$\subsysA$ 端在 $z$ 基中的约化密度矩阵。$\varphi$ 是联合态第二项相对第一项的相位。取 $0$ 或 $\pi$ 时，对角元相同，两个非对角元的符号相反。}
        \label{fig:qm-reduced-spin-density-matrices}
    \end{notecaptionblock}
    若把原态两项之间的加号改为减号，两端都测 $z$ 时的联合概率与原来相同。此时 $\op\rho_{\subsysA}$ 的两个非对角元反号（图\ref{fig:qm-reduced-spin-density-matrices}），$p_{\subsysA}(+x)$ 变为 $(1 - 1/\sqrt2)/2$。只知道沿 $z$ 的计数，就无法确定转向 $x$ 后的测量概率。
\end{example}

能否分别给两端各指定一个纯态，从而完整描述联合制备？
若 $\ket{\Psi} = \ket{u}_{\subsysA} \otimes \ket{v}_{\subsysB}$，系数矩阵 $C$ 的每一行
都是同一行向量$(v_1,v_2,\ldots)$的倍数。上例的两行并不成比例，故这个态无法写成乘积，
称为纠缠的纯态。一般情形可用下面的方法找出所需的乘积项数。

\begin{theorem}{二分纯态的Schmidt分解}{qm-schmidt}
    对有限维的两个系统，归一化联合纯态总可写成 Schmidt 分解
    \begin{equation}
        \ket{\Psi} = \sum_r\sqrt{\lambda_r}\,\ket{a_r}_{\subsysA} \otimes \ket{b_r}_{\subsysB},
        \qquad \lambda_r > 0,\quad\sum_r\lambda_r = 1.
    \end{equation}
    两组局部态各自正交归一，$\lambda_r$ 是两端约化态共有的非零本征值。
\end{theorem}
\begin{proof}
    选 $\op{\rho}_{\subsysA}$ 的正交本征基 $\ket{a_r}$，相应本征值为 $\lambda_r$。
    将 $\subsysA$ 投影到 $\ket{a_r}$，$\subsysB$ 留下的未归一化态记为 $\ket{w_r}$。由部分迹的分量公式，
    \begin{subequations}
        \begin{align}
            \ket{w_r} &= (\bra{a_r} \otimes \one_{\subsysB})\ket{\Psi},\\
            \braket{w_s}{w_r} &= \mel{a_r}{\op{\rho}_{\subsysA}}{a_s} = \lambda_r\delta_{rs}.
        \end{align}
    \end{subequations}
    因此非零 $\ket{w_r}$ 彼此正交，将每项除以 $\sqrt{\lambda_r}$ 得到 $\ket{b_r}$。
    把 $\ket{\Psi}$ 沿 $\subsysA$ 的这组基展开就得到上式。分别对一端取迹，两端都留下以 $\lambda_r$ 为权重的正交投影之和。
\end{proof}

\begin{proof}[用系数矩阵计算Schmidt分解]
    取前面乘积基中的系数矩阵 $C$，作奇异值分解(singular value decomposition, SVD) $C = UDV\adjoint$。
    $U,V$的列正交归一，$D$的非零对角元记为$\sqrt{\lambda_r}$。
    定义
    \begin{subequations}
        \begin{alignat}{1}
            \ket{a_r} &= \sum_aU_{ar}\ket{a},\\
            \ket{b_r} &= \sum_bV_{br}\conj\ket{b}.
        \end{alignat}
    \end{subequations}
    代回乘积基展开，奇异值分解的每个非零奇异值就对应一个 Schmidt 项。
    对 $\subsysB$ 取部分迹，利用 $\braket{b_s}{b_r} = \delta_{sr}$，得到
    \begin{equation}\label{eq:qm-schmidt-reduced-states}
        \op{\rho}_{\subsysA} = \sum_r\lambda_r\ketbra{a_r}{a_r}.
    \end{equation}
\end{proof}
\begin{corollary}[由一端的纯度判别二分纯态的纠缠]
    对联合纯态，只要一端的约化态是纯态，联合态就是乘积态，反之亦然。
    因为一端的约化态为纯态恰好表示只有一个非零本征值，而此时 Schmidt 分解只含一个乘积项。
\end{corollary}
约化态(reduced state)的 von Neumann 熵量化纯联合态的二分纠缠。$S$ 将密度算符映为一个数。计算时先取标量函数 $s(\lambda) = -\lambda\ln\lambda$，它在 $\lambda = 0$ 处按连续极限取 $s(0) = 0$，再对 $s(\op\rho_{\subsysA})$ 取迹，
\begin{equation}\label{eq:qm-entanglement-entropy}
    S(\op{\rho}_{\subsysA}) = -\sum_r\lambda_r\ln\lambda_r .
\end{equation}
当联合态是混合态时，局部熵也包含随机制备所贡献的不确定性。
\begin{definition}{可分态与纠缠}{qm-separability}
    两个有限维系统的联合态若可由下式制备，则称为可分态(separable state)，
    \begin{equation}\label{eq:qm-separable-state}
        \op{\rho}_{\subsysAB} = \sum_r p_r\,\op{\rho}_{\subsysA}^{(r)} \otimes \op{\rho}_{\subsysB}^{(r)},\qquad
        p_r \geq 0,\qquad \sum_r p_r = 1 .
    \end{equation}
    即先以概率 $p_r$ 选择共同标签 $r$，再按标签分别制备两端。
    无法这样制备的联合态称为纠缠态(entangled state)。对纯态，可分条件化为态矢本身是乘积态。
\end{definition}

\begin{example}[unbreakable]{自旋单态的双轴关联}{qm-singlet-density}
    两个自旋为 $1/2$ 的粒子处于单态 $\ket{\singlet} = (\ket{\up,\dn} - \ket{\dn,\up})/\sqrt2$，箭头给出沿 $z$ 轴的自旋投影方向。求第一粒子的约化密度算符。若两端各沿单位向量 $\bm a,\bm b$ 测量自旋，并把每端读数记为 $\pm1$，求两读数乘积的期望值。
    \par\smallskip\textbf{解}\quad
    由 $\ketbra{\singlet}{\singlet}$ 写出单态密度算符的 Pauli 展开，
    \begin{equation}
        \op{\rho}_{\singlet} = \frac14\Bigp{\one \otimes \one
            - \sum_{i=x,y,z}\op{\sigma}_i \otimes \op{\sigma}_i}.
    \end{equation}
    对第二个自旋取迹，
    $\tr\one = 2$、$\tr\op{\sigma}_i = 0$给出$\op{\rho}_{\subsysA} = \one/2$。
    但联合相关为
    \begin{equation}
        \begin{alignedat}{1}
            \ave{\op{\sigma}_i \otimes \op{\sigma}_j} &= \tr\bigb{\op{\rho}_{\singlet}(\op{\sigma}_i \otimes \op{\sigma}_j)}\\
            &= -\frac14\sum_k
                \tr(\op{\sigma}_k\op{\sigma}_i)\tr(\op{\sigma}_k\op{\sigma}_j)
                = -\delta_{ij}.
        \end{alignedat}
    \end{equation}
    将两端的测量轴分别取在单位向量 $\bm a,\bm b$ 方向，对分量求和便得到自旋单态的关联
    \begin{equation}
        E(\bm a,\bm b)
        = \bigave{(\bm a \cdot \opbm{\sigma}) \otimes (\bm b \cdot \opbm{\sigma})}
        = -\bm a \cdot \bm b.
        \label{eq:singlet-correlation}
    \end{equation}
\end{example}

\subsectionstar{两端各自测量能否确定联合态}
\label{ss:qm-local-tomography}
在\secref{sec:interference}的单路相移中，复数可用一对实坐标表示。现在考虑两台分开的自旋分析器。
对许多次相同的制备，两端各自选择测量轴，再逐次配对读数。把各轴组合都测遍，能否由这些局部测量的联合统计重建完整的密度算符？
这种能力称为局部层析（local tomography）\cite{hardy-wootters}。对两个复二维系统，令 $\op\sigma_0 = \one$。
十六个乘积 $\op\sigma_i \otimes \op\sigma_j$（$i,j = 0,x,y,z$）构成联合系统自伴算符空间的一组正交基。
每个乘积的期望都能从两端配对的读数求得，下标为零的一端只需记常数 $1$。这些期望就给出了重建联合态所需的全部系数。

若改选每端为二维实 Hilbert 空间，并用通常的实张量积合成两端，单端可观测量在 $z$ 基中都由实对称矩阵表示。这些矩阵由单位矩阵和式\eqref{eq:qm-pauli-matrices}中的 $\sigma_x,\sigma_z$ 张成。
同式中的 $\sigma_y$ 是纯虚矩阵，$\sigma_y \otimes \sigma_y$ 却是实对称矩阵，在联合乘积基中表示可观测量 $\op M$。它满足 $\op M^2 = \one$、$\tr\op M = 0$。
因此 $\op\rho_\pm = (\one \pm \op M)/4$ 是两个合法的联合态，整体测量 $\op M$ 时分别必得 $+1$ 与 $-1$。
但单端可观测量 $\op A,\op B$ 的矩阵 $A,B$ 都对称，而 $\sigma_y$ 反对称，故 $\tr(A\sigma_y) = \tr(B\sigma_y) = 0$，
所以 $\tr[(\op\rho_+ - \op\rho_-)(\op A \otimes \op B)] = 0$。两端各自测量能得到的联合概率全都相同。
若要求局部层析，通常的实张量积在这个例子中就无法满足要求。在复数理论中，两端只要分别测 $y$ 分量并比较读数，便能区分这两个态。

\section{路径标记、量子擦除与暗端计数}
\label{sec:qm-path-readout}
沿用图\ref{fig:two-path-interferometer}的平衡干涉仪。每次送入一个光子，先使两臂都畅通，调节相移器直到光子全部从正端出射，把这个位置记作 $\varphi = 0$。
保持分束器、合束器和两端的探测器固定。先在臂中给路径加上可读的标记，再改变标记的读出基，最后撤去标记、在一臂放入吸收体。三种操作会怎样改变同一对探测器的计数？

\subsection{能从记录辨认哪条路径}
\label{ss:qm-path-marking}
可让 $a$ 臂中的光子保持水平偏振，在 $b$ 臂放入波片，转动其偏振方向。让两个输出端接收所有偏振。
此时偏振充当记录器，它随光子一起到达合束器。即使无人读偏振，条纹可见度也会改变。

\Needspace{4\baselineskip}
用 $\ket{D_a},\ket{D_b}$ 表示两臂留下的归一化记录态。按“路径、记录”的顺序写成 $\ket{a,D_a} = \ket a \otimes \ket{D_a}$，则耦合后的联合态为
\begin{equation}
    \ket{\Psi} = \frac{\ket{a,D_a} + \e^{\i\varphi}\ket{b,D_b}}{\sqrt2}.
\end{equation}

联合密度算符展开含四项，其中一个交叉项的算符部分可以按子系统分开，
\[
    \ketbra{a,D_a}{b,D_b} = (\ketbra{a}{b}) \otimes (\ketbra{D_a}{D_b}).
\]
对记录器取迹，第二个因子给出内积 $\braket{D_b}{D_a}$。因此
\begin{equation}
    \op{\rho}_{\pathstate} = \frac12\Bigb{
        \ketbra{a}{a} + \ketbra{b}{b}
        + \e^{-\i\varphi}\braket{D_b}{D_a}\ketbra{a}{b}
        + \e^{\i\varphi}\braket{D_a}{D_b}\ketbra{b}{a}}.
\end{equation}

对记录器取迹后，路径约化态的 $ab$ 矩阵元带有因子 $\braket{D_b}{D_a}$。
在正输出端测得 $\ket{+}$ 的概率为
\begin{equation}
    P_+ = \frac12\bigb{1 + \Re\bigp{\e^{\i\varphi}\braket{D_a}{D_b}}}.
\end{equation}
扫描相位 $\varphi$，记 $P_+$ 的最大、最小值为 $P_{\max},P_{\min}$。在等路径权重的理想模型中，条纹可见度(visibility)为
\begin{equation}
    \mathcal V = \frac{P_{\max} - P_{\min}}{P_{\max} + P_{\min}}
              = \abs{\braket{D_a}{D_b}} .
\end{equation}
记录态相同时，两路都有相同的偏振，通过正端的概率随相位在零和一之间变化。若两臂的偏振互相正交，测偏振就能可靠地区分路径（which-way measurement），只数正端光子的探测器则对所有相位都给出一半的通过率。

\phantomsection\label{exp:qm-path-1998}1998 年，Dürr 等用驻波光栅分开铷原子的物质波，在远处测得的原子位置分布中看到了干涉条纹。他们用微波脉冲给两条路径标上互相正交的原子内部态，原子位置分布中的条纹就消失了。实验估算出微波给原子的动量太小，无法把条纹抹平。即使探测时只记录原子的位置，两个内部态的内积为零也足以消去干涉项\cite{durr-whichway}。

\phantomsection\label{exp:qm-recoil-2025}2025 年，Zhang 等用光镊束缚单个铷原子，将它冷却到接近运动基态。他们沿同一轴线从两侧收集散射光子。光子向两侧散射时，原子沿这条轴线得到的反冲动量方向相反。将两路光引到分束器合束，扫描相对相位，就能测出单光子计数的起伏。
减小光镊的阱深，原子的初始动量分布变窄，两种反冲后的运动态重叠得更少，所测光子条纹的可见度也随之降低。原子的动量越能区分光子的散射方向，光子计数所显示的干涉就越弱\cite{recoiling-atom}。

\subsection{量子擦除与条件条纹}
\label{ss:qm-quantum-eraser}
回到偏振标记装置。在输出端再加偏振分析器，把每次的路径端口读数与偏振读数配对保存。

设 $\ket{D_a}$ 与 $\ket{D_b}$ 正交，改在记录器的基 $\ket{D_\pm} = (\ket{D_a} \pm \ket{D_b})/\sqrt2$ 中测量，按读数把粒子分成两组。
沿 $\ket{D_\pm}$ 取联合态的分量，得到 $(\ket a \pm \e^{\i\varphi}\ket b)/2$，其范数平方为 $1/2$，即各组出现的概率。
组内归一化后，路径态为 $(\ket a \pm \e^{\i\varphi}\ket b)/\sqrt2$。在合束器的正端计数，两组的通过率和合并后的通过率分别为
\begin{subequations}\label{eq:qm-conditional-interference}
    \begin{align}
        \P(+ \mid D_\pm) &= \frac{1 \pm \cos\varphi}{2},\\
        P_+ &= \frac12\P(+ \mid D_+) + \frac12\P(+ \mid D_-) = \frac12 .
    \end{align}
\end{subequations}
图\ref{fig:qm-quantum-eraser-fringes}中，一组条纹达到峰值时，另一组恰好达到谷值。两组各占一半，合并后的正端概率始终为 $1/2$。这种按另一组基分析路径标记、在条件样本中读出干涉的操作称为量子擦除。

\begin{figure}[htbp]
    \centering
    \tikzfigure{quantum-eraser-fringes}
    \caption{扫描两臂的相对相位 $\varphi$，比较合束器的正端概率。实线和虚线分别以偏振读数 $D_+$ 和 $D_-$ 的组内事件总数为分母，两组条纹相差半个周期。点线以两组的事件总数为分母，给出合并后的正端概率 $P_+ = 1/2$。}
    \label{fig:qm-quantum-eraser-fringes}
\end{figure}

若两臂原来的标记分别为水平和竖直偏振，这一步用的就是与水平轴成 $\pm45^\circ$ 的两个线偏振出口。已经保存的正、负端口总数在分组前后相同。

\phantomsection\label{exp:qm-eraser-2002}2002 年，Walborn 等将一对偏振纠缠光子中的一个送向双缝，另一个送到偏振分析器。两缝前各有一片四分之一波片，快轴相互垂直。两条路径对应正交的双光子偏振态。要靠偏振认出经过哪条缝，须比较伙伴光子的线偏振和过缝光子的圆偏振。为观察擦除后的条纹，研究者移动双缝后的探测器，在各个位置统计它与伙伴光子的符合计数。伙伴端的偏振片分别沿两片波片的快轴取向时，测到的两组条纹峰谷互换。合并两种偏振的结果，干涉条纹消失\cite{eraser-double-slit}。

\subsection{好炸弹与坏炸弹的暗端检验}
\label{ss:qm-absorber-test}
Elitzur 与 Vaidman 的思想实验可以在同一台干涉仪上用光学吸收体说明\cite{ev-object,kwiat-object}。
撤去偏振标记和分析器，将相位固定在原来的零点。待检物放入 $b$ 臂。“好炸弹”遇到一个光子就触发，模型取为理想吸收体；
“坏炸弹”让光子无损通过，保持原相位。后一种情形两路仍在负端相消，所以负端称为暗端。

物体能吸收时，第一个分束器把一半概率送往吸收通道。在没有吸收的试验中，到达合束器的未归一化态为 $\ket a/\sqrt2$。
其正、负端振幅分别为 $\braket{+}{a}/\sqrt2$、$\braket{-}{a}/\sqrt2$，两者都等于 $1/2$。以全部入射光子为分母，就得到\tabref{tab:qm-dark-port}的三种计数概率。

\begin{table}[H]
\centering
\caption{干涉仪调至零相位时，每个入射光子的三种互斥结果。}
\label{tab:qm-dark-port}
\begin{tabular}{lccc}
    \toprule
    $b$ 臂中的物体 & 被吸收 & 正端点击 & 暗端点击 \\
    \midrule
    透明且保持相位 & $0$ & $1$ & $0$ \\
    理想吸收体 & $1/2$ & $1/4$ & $1/4$ \\
    \bottomrule
\end{tabular}
\end{table}

在这两个备选假设下，暗端的一次点击就确认了吸收体存在，而且这次的光子到达了探测器。正端点击与两种假设都相容。
对一个理想吸收体，一次检验既辨认出它又保留它的概率为 $1/4$。若只在光子幸存的试次中计数，一半会到达暗端。

\phantomsection\label{exp:qm-absorber-1995}1995 年，Kwiat 等将一对光子中的一个送往触发探测器，另一个送入 Michelson 干涉仪。两臂畅通时，先把一个出口的符合计数调到最低，测出暗端背景。再将一面可移动的镜子插入一臂，把这一路的光引向吸收探测器，暗端与触发端的符合计数随即升高。暗端收到的光子来自畅通的一臂，它能从这里出射，是因为另一路受阻后两路不再在暗端相消\cite{kwiat-object}。

\section{Bell实验与CHSH不等式}
单端统计需要对另一端所有可能的结果求和。若保留两端的配对记录，就能检验关联从何而来。
同一方向上两个自旋总给出相反读数，也可以由源预先给两端分发相反的指令解释。
Bell实验让两端各自改变测量方向，比较四组关联能否同时由一套局域指令给出。

\subsection{实验设置与关联估计}
\label{ss:bell-counting}

Bell 实验要把每对粒子的两端读数配成一次联合计数，并按两端的设置组合分别统计（图\ref{fig:bell-pair}）。

\begin{figure}[htbp]
    \centering
    \tikzfigure{bell-pair}
    \caption{从中央的源 $\mathrm{S}$ 沿箭头向两端读：方框内的 $a,a'$ 与 $b,b'$ 分别是 Alice 和 Bob 可选的方向；每端在所选方向的两个出口之一记录 $+1$ 或 $-1$。}
    \label{fig:bell-pair}
\end{figure}

对每组设置$(a,b)$，分别记录同号计数$N_{++},N_{--}$
和异号计数$N_{+-},N_{-+}$。
若每次试验两端都登记为 $\pm1$，这四项计数之和就是总次数 $N_{ab}$。关联的实验估计为
\begin{equation}
    E_{\mathrm{obs}}(a,b) =
    \frac{N_{++} + N_{--} - N_{+-} - N_{-+}}{N_{ab}}.
\end{equation}
同号结果对乘积贡献$+1$，异号结果贡献$-1$，
所以这个式子正是两端结果乘积的样本平均。
有漏检时，可以事先为未读出的结果指定一个符号，把这次试验也计入。若只统计两端都读出的试验，所得是符合样本的条件平均，直接套用下文的界还需额外的抽样假设。

要用四组关联检验同一个局域模型，需要假定每对粒子携带的信息与设置选择独立。为检验局域性，可以将两端分开，各自在测量前随机选轴，使一端的选择即使以光速传播也来不及影响另一端的读数\cite{belltests}。

\subsection{局域模型的界与量子预测}
\label{ss:chsh-derivation}

\begin{theorem}{CHSH不等式}{qm-chsh}
    设两端各选两种设置之一，每次分别得到 $\pm1$。若给定共同隐藏变量 $\lambda$ 后，
    两端联合概率分解为各自局部条件概率的乘积 $\P(A \mid a,\lambda)\P(B \mid b,\lambda)$，且 $\lambda$ 的分布与设置选择独立，
    则四组结果乘积的平均值 $E$ 满足
    \begin{equation}\label{eq:qm-chsh-local-bound}
        \abs{S} = \abs{E(a,b) + E(a,b') + E(a',b) - E(a',b')} \leq 2 .
    \end{equation}
\end{theorem}
\begin{proof}
    在局域隐藏变量模型(local hidden variable model)中，可把两端的局部随机数一起记入 $\lambda$。
    固定 $\lambda$ 后，左端的结果 $A(a,\lambda)$ 只由设置 $a$ 决定，右端的结果 $B(b,\lambda)$ 只由设置 $b$ 决定。
    模型因此给每一种本地设置都指定了一个结果，即使当次实验只在每端选了一种设置。暂略去共同的 $\lambda$，四个结果 $A(a),A(a'),B(b),B(b')$ 各取 $\pm1$，所以
    \begin{equation}
        A(a)[B(b) + B(b')] + A(a')[B(b) - B(b')] = \pm2.
    \end{equation}
    若 $B(b) = B(b')$，两数之和为 $\pm2$，差为零。若 $B(b) = -B(b')$，和为零，差为 $\pm2$。
    无论 Alice 的两个结果如何取值，四项组合的绝对值都是 2。记这个组合为 $C(\lambda)$。
    设隐藏变量的归一化概率密度为 $p(\lambda)$。在实际选定的一组设置 $(a,b)$ 下，关联函数(correlation function)是两端结果乘积的平均值，
    \begin{equation}
        E(a,b) = \int\d\lambda\,p(\lambda)A(a,\lambda)B(b,\lambda).
    \end{equation}
    设置与$\lambda$分布独立，使四个关联都能用同一个$p(\lambda)$平均。
    由三角不等式，平均后的 CHSH 组合满足
    \[
        \abs{S} = \biggabs{\int p(\lambda)C(\lambda)\d\lambda}
        \leq \int p(\lambda)\abs{C(\lambda)}\d\lambda = 2 .
    \]
\end{proof}
自旋单态 $\ket{\singlet}$ 的量子预测由式\eqref{eq:singlet-correlation}给出。取 Alice 的测量方向为 $\bm a = \bm e_z$、$\bm a' = \bm e_x$，
Bob 的测量方向为 $\bm b = (\bm e_z + \bm e_x)/\sqrt2$、$\bm b' = (\bm e_z - \bm e_x)/\sqrt2$。四个点积分别给出
\begin{subequations}
    \begin{alignat}{2}
        E(a,b) &= -\frac1{\sqrt2},
        &\quad E(a,b') &= -\frac1{\sqrt2},\\
        E(a',b) &= -\frac1{\sqrt2},
        &\quad E(a',b') &= \frac1{\sqrt2}.
    \end{alignat}
\end{subequations}
于是$S = -2\sqrt2$，其绝对值超过式\eqref{eq:qm-chsh-local-bound}给出的局域界 $2$。
这里使用的是自旋测量轴；线偏振光子实验的物理偏振角与Bloch球角相差二倍，
选取光学分析器角度时应转换这一约定。

\phantomsection\label{exp:qm-bell-1972}1972 年，Freedman 与 Clauser 利用钙原子级联跃迁发出光子对。两端各放一块偏振片，改变两片的夹角，测量两端同时探测到光子的速率。不同夹角下的符合率违反了他们检验的 Bell 不等式。由于只有部分光子能被探测到，推导该实验的局域界时还假定，对已经到达探测器的光子，记录效率与它此前是否通过偏振片无关\cite{belltests}。

\phantomsection\label{exp:qm-bell-1982}1982 年，Aspect 等在两端各用声光开关把飞行中的光子导向两块取向不同的偏振片。切换方向所用的时间短于光在两端之间的飞行时间，两端的偏振关联仍违反 Bell 不等式。这次实验采用预设的频率切换，设置随时间的变化可以预知\cite{belltests}。

\Needspace{5\baselineskip}
\phantomsection\label{exp:qm-bell-2015}2015 年，Hensen 等在相距 $\qty{1.3}{km}$ 的两处分别测量金刚石中的电子自旋。中间站探测来自两地的光子，用特定的符合信号标记纠缠制备成功的试次。两端各自随机选轴并迅速读出自旋。一端选轴的信息即使以光速传播，也来不及影响另一端当次的读数。每个有制备成功信号的试次都计入统计，无须假定检出的样本能代表其余试次。245 次试验给出 $S_{\mathrm{obs}} = 2.42 \pm 0.20$。式\eqref{eq:qm-chsh-local-bound}限制的是理论期望，有限样本会有涨落。按论文的统计检验，在局域模型下，观测到同样大或更大偏离的概率至多为 $0.039$\cite{belltests}。

若 Bob 操作后保留全部测量结果，Alice 能否仅从本地计数察觉这一操作？
设 $\op{M}_b$ 只作用在 Bob 的空间，$b$ 标记一个可能的结果。对结果求和，联合态成为
\begin{gather*}
    \op{\rho}' = \sum_b(\one_A \otimes \op{M}_b)\op{\rho}(\one_A \otimes \op{M}_b\adjoint),\\
    \sum_b\op{M}_b\adjoint \op{M}_b = \one_B .
\end{gather*}
右侧的完备条件保证概率总和为一。

现在计算 Alice 端任意可观测量 $\op{O}_A$ 在操作后的期望值。由迹的循环性质，
\begin{equation}
    \begin{alignedat}{1}
        \ave{\op{O}_A}' &= \sum_b
            \tr\Bigb{\op{\rho}\,
            (\one_A \otimes \op{M}_b\adjoint)(\op{O}_A \otimes \one_B)
            (\one_A \otimes \op{M}_b)}\\
        &= \tr\Bigb{\op{\rho}\,
            \bigp{\op{O}_A \otimes \sum_b\op{M}_b\adjoint \op{M}_b}}
            = \tr\bigb{\op{\rho}(\op{O}_A \otimes \one_B)}.
    \end{alignedat}
\end{equation}
所有局部平均值相同，故 $\op{\rho}_A$ 不变，Bob 无法借此向 Alice 传递信号。Alice 收到 Bob 发来的结果标签后，才可将样本分组并比较各组的条件态。

\section{一般量子测量}
\label{sec:povm}
Bell 关联取决于分析器记录的正负号。如果读出电路偶尔将正、负标签互换，怎样给出所记录的概率？按这个标签筛出的粒子，又处于什么态？

设每次记录独立地以概率 $\epsilon$ 发生标记错误，$\op\Pi_+$、$\op\Pi_-$ 分别投影到正、负出口所对应的自旋态。正记录可能来自正确标记的正出口，也可能来自误记的负出口。
用式\eqref{eq:density-statistics}计算实际出口的概率，再按标记正确或错误的概率加权，得到
\begin{subequations}
    \begin{align}
        p_+ &= \tr(\op{\rho}\op E_+),\\
        \op E_+ &= (1 - \epsilon)\op\Pi_+ + \epsilon\op\Pi_- .
    \end{align}
\end{subequations}
若 $0 < \epsilon < 1/2$，其本征值 $1 - \epsilon$ 和 $\epsilon$ 均在 0 与 1 之间，故 $\op{E}_+$ 已超出投影算符的范围。

\begin{definition}{离散测量的POVM与仪器映射}{qm-povm-instrument}
    仪器映射(instrument map)将输入态 $\op{\rho}$ 映到记录 $m$ 对应的未归一化输出态，
    \begin{equation}\label{eq:qm-measurement-instrument}
        \mathcal I_m(\op{\rho}) = \sum_\alpha \op{M}_{m\alpha}\op{\rho} \op{M}_{m\alpha}\adjoint .
    \end{equation}
    其中 $\alpha$ 标记同一记录下未进一步分辨的作用通道。相应的效应算符(effect operator)满足
    \begin{subequations}
        \begin{align}
            \op{E}_m &= \sum_\alpha \op{M}_{m\alpha}\adjoint \op{M}_{m\alpha} \geq 0,\\
            \sum_m\op{E}_m &= \one,
        \end{align}
    \end{subequations}
    这样一族正算符 $\{\op{E}_m\}$ 称为正算符值测度(positive operator-valued measure, POVM)。结果概率及条件态为
    \begin{subequations}
        \begin{align}
            p_m &= \tr(\op{E}_m\op{\rho}),\\
            \op{\rho}_m &= \frac{\mathcal I_m(\op{\rho})}{p_m}
            \quad(p_m > 0).
        \end{align}
    \end{subequations}
\end{definition}
由迹的循环性质，$\tr\mathcal I_m(\op{\rho}) = \tr(\op{E}_m\op{\rho}) = p_m$，这解释了条件态的归一化因子。各结果相加，由完备条件得到总迹 $\tr\op{\rho}$。
在误记模型中，正记录的两个测量算符分别为 $\op{M}_{+,1} = \sqrt{1 - \epsilon}\,\op\Pi_+$、$\op{M}_{+,2} = \sqrt{\epsilon}\,\op\Pi_-$。
两项的输出各保留实际出口的自旋，输出的迹相加便是 $p_+$。

理想 Lüders 测量取 $\op{M}_m = \op\Pi_m$。若在出口 $m$ 接上幺正变换 $\op{U}_m$，新算符 $\op{U}_m\op\Pi_m$ 给出相同的概率，却可改变测后态。

\begin{example}{单态与混合态的鉴别}{qm-singlet-mixture-test}
    一台装置制备双自旋单态，另一台以相等概率制备 $\ket{\up,\dn}$ 或 $\ket{\dn,\up}$，箭头均沿 $z$ 轴。两台装置在每一端的局部态都是 $\one/2$。设计一种联合测量，说明怎样从重复实验的计数分布区分这两种制备。
    \par\smallskip\textbf{解}\quad
    分别沿 $x$ 轴测量两端自旋，统计同号结果的比例。单态在同一方向上完全反相关，同号概率为零。混合态的每一个乘积态在两端各以相等概率产生 $+x$、$-x$，同号概率为 $1/2$。
\end{example}
\section{量子查询算法与 Deutsch 问题}
\label{sec:qm-deutsch-query}
干涉仪的正、负端可以读出两路振幅的相对符号。Deutsch 问题把这个读出用来判断一个未知函数的性质。
给定一个只能调用的装置（黑箱），其中存有函数 $f$，输入 $0,1$ 各对应一个取值为 $0$ 或 $1$ 的输出。
要求准确判断 $f(0)$ 和 $f(1)$ 是否相同。若每次送入一个经典比特并读取相应函数值，
一次查询只能知道其中一个函数值。另一个值既可能与它相同，也可能与它不同，因此要保证答对就必须查询两次。

一个量子比特（qubit）是一套有两个正交基态 $\ket 0,\ket 1$ 的量子系统。量子电路中称可控的幺正操作为量子门。
本题允许相干地调用黑箱，两比特查询门的作用由下面的规则及其线性延拓给定\cite{cleve-algorithms}，
\begin{equation}\label{eq:qm-deutsch-oracle}
    \op U_f\ket{x,y} = \ket{x,y\mathbin{\oplus}f(x)},\qquad x,y\in\{0,1\}.
\end{equation}
这里 $\oplus$ 为模二加法。固定 $x$，若 $f(x) = 1$，第二比特的两个基态就互换。
制备 $\ket{\pm} = (\ket0 \pm \ket1)/\sqrt2$，令第一比特处于 $\ket+$，第二比特处于 $\ket-$。
互换 $\ket-$ 的两个分量会给它整体乘上 $-1$。因而一次查询后的联合态是
\[
    \op U_f(\ket+\otimes\ket-)
       = \frac{(-1)^{f(0)}\ket0 + (-1)^{f(1)}\ket1}{\sqrt2}\otimes\ket- .
\]
查询后第二比特仍处于 $\ket-$。若 $f(0) = f(1)$，第一比特的两个振幅同号，所得态与 $\ket+$ 只差一个整体相位。若 $f(0) \ne f(1)$，两振幅异号，第一比特就处于 $\ket-$（略去整体相位）。
以 $\ket\pm$ 为基读取第一比特就能确定答案。实际若沿 $\ket0,\ket1$ 读出，可先施加 Hadamard 门，
它将 $\ket+$ 送到 $\ket0$，将 $\ket-$ 送到 $\ket1$。末端的一位记录给出两函数值的奇偶性 $f(0)\oplus f(1)$。
这里比较的是同一个查询任务所需的黑箱调用次数，搭建查询门和保持相干所需的资源须另行计算。

\reading{Bell 实验见文献\cite{belltests}，部分迹与量子操作见文献\cite{manzano}。量子比特在有噪声时如何存储，见\secref{sec:qm-error-correction}。下一章从远处的散射计数反推相互作用。}
